% Agent 09: Pages 377--380 (PDF pp.18--20)
% Section 2.6 Radiative Transfer (end): optical depth, optically thick/thin,
%   diffusion approximation, Rosseland mean opacity
% Section 2.7 Shock Waves (beginning): types of shocks, relativistic
%   Rankine-Hugoniot equations, Figure 2.7.1
% Equations (2.6.26)--(2.6.47), (2.7.1a)--(2.7.1d)
% =============================================================================

%%% ---------------------------------------------------------------
%%% Continuation of S2.6  Radiative Transfer
%%% (picking up from the equation of radiative transfer, eq.~(2.6.22))
%%% ---------------------------------------------------------------

Notice that for matter in thermodynamic equilibrium the relationship
$\varepsilon_\nu / 4\pi\kappa_\nu = B_\nu$ permits one to rewrite the equation
of radiative transfer~(2.6.22) in the form
%
\begin{equation}
\tag{2.6.26}
\frac{dI_\nu}{dl} + (3n\cdot a + \dot\theta + 3\sigma^{\alpha\beta}
n_\alpha n_\beta)\, I_\nu
= \rho_0\,\kappa_\nu\bigl(B_\nu - I_\nu\bigr)
  + \left.\frac{dI_\nu}{dl}\right|_{\text{scattering}}
  + \left.\frac{d(I_\nu/\nu^3)}{d\tau_\nu}\right|_{\text{scattering}} ;
\end{equation}
%
or, equivalently,
%
\begin{equation}
\tag{2.6.27}
\frac{d(I_\nu/\nu^3)}{d\tau_\nu}
= \rho_0\,\kappa_\nu \left(\frac{B_\nu}{\nu^3}
  - \frac{I_\nu}{\nu^3}\right)
  + \left.\frac{d(I_\nu/\nu^3)}{d\tau_\nu}\right|_{\text{scattering}}.
\end{equation}

%%% ---------------------------------------------------------------
%%% Optical depth
%%% ---------------------------------------------------------------

\bigskip
\noindent
\emph{Optical depth}

\medskip
\noindent
Consider radiation propagating out of a medium into surrounding empty space.
Follow a given ray ``backward'', from ``infinity'' into the medium.  Let the ray
have frequency $\nu_\infty$ at infinity; then its frequency at location~$l$
($l = {}$proper distance measured in LRF of medium) is
%
\begin{equation}
\tag{2.6.28}
\nu(l) = \nu_\infty \exp\!\left[\,-\int_i^l
\bigl(\dot\theta + n^\alpha n^\beta \sigma_{\alpha\beta}\bigr)\,dl\right].
\end{equation}
%
[cf.\ eq.~(2.6.14)].  In calculating the change of intensity along the ray, it is
often useful to replace the proper-length parameter~$l$ by the
``optical-depth'' parameter
%
\begin{equation}
\tag{2.6.29}
\tau_\nu = \int_i^l \rho_0\,\kappa_\nu\,dl.
\end{equation}

In the integration $\kappa_\nu$ must be evaluated at the frequency
$\nu(l)$.  The optical depth $\tau_\nu$ depends on (i) the world line of
the chosen ray, (ii) the frequency of the ray in spacetime, (iii) location
along that ray.  Then the equation of transfer
world line; and (iii) location along the ray at ``infinity'' $\nu_\infty$\ldots

One can also introduce an optical depth for scattering radiation out of the
beam, $\tau_s$:
%
\begin{equation}
\tag{2.6.30}
\tau_s = \int_i^l \rho_0\,\kappa_s\,dl.
\end{equation}

In terms of optical depths, the law of radiative transfer~(2.6.27) reads
%
\begin{equation}
\tag{2.6.31}
\frac{d(I_\nu/\nu^3)}{d\tau_\nu}
= \frac{B_\nu}{\nu^3} - \frac{I_\nu}{\nu^3}
  - \frac{\kappa_s}{\kappa_\nu}
  \left.\frac{d(I_\nu/\nu^3)}{d\tau_\nu}\right|_{\text{scattering}},
\end{equation}
%
where
%
\begin{equation}
\tag{2.6.32}
\left.\frac{d(I_\nu/\nu^3)}{d\tau_\nu}\right|_{\text{scattering}}
= \frac{\kappa_s}{\kappa_\nu}\left[\frac{I_\nu}{\nu^3} - \frac{1}{\kappa_s}
  \int \frac{dk_s}{d\Omega\,d\nu'}\,
  (\mathbf{p}',\mathbf{p})\,I_\nu'\,d\Omega'\,d\nu'
  - \rho_0\,\kappa_s\,I_\nu \right].
\end{equation}

%%% ---------------------------------------------------------------
%%% Optically thin / optically thick
%%% ---------------------------------------------------------------

One says that the medium is \emph{optically thin} if everywhere along
the ray $\tau_\nu \ll 1$ (or $\tau_s \ll 1$).

Consider a medium that (i) has its matter in thermodynamic equilibrium with
a spatially uniform temperature, (ii) is optically thick along a chosen ray, and
(iii) has absorption dominant over scattering, i.e.,
%
\begin{equation}
\tag{2.6.33}
\kappa_\nu \gg \kappa_s.
\end{equation}

One says that a medium is \emph{optically thick} to emission and absorption (or to
scattering) along a given ray if optical depths $\tau_\nu \gg 1$
(or $\tau_s \gg 1$) are achieved.

\medskip
\noindent
Along that ray, then the radiation emerging to infinity along the chosen ray must
have the blackbody form
%
\begin{equation}
\tag{2.6.34}
\frac{I_\nu}{\nu^3} = \frac{B_\nu}{\nu^3}
= \frac{2h\nu^3/c^3}{e^{h\nu/kT} - 1}.
\end{equation}
%
[One can prove this easily from the equation of transfer~(2.6.31).]\
Moreover, at a point in the medium where all rays are optically thick, the
radiation will have a blackbody intensity at all frequencies and in all
directions; so the energy density and pressure in the radiation will be
%
\begin{equation}
\tag{2.6.35}
\rho^{\text{rad}} = 3p^{\text{rad}} = bT^4,
\end{equation}
%
where $b$ is the universal constant
%
\begin{equation}
\tag{2.6.36}
b = \frac{8\pi^5 k^4}{15c^3 h^3}
  = 7.56 \times 10^{-15}\;\frac{\text{ergs}}{\text{cm}^3\,\text{K}^4}.
\end{equation}

Consider, alternatively, a medium that is optically thin along a chosen ray,
and has negligible scattering ($\tau_s \ll 1$) along that ray.  Then the
equation of transfer~(2.6.31) predicts for the radiation emerging to infinity
%
\begin{equation}
\tag{2.6.37}
I_\nu\,\nu^3 = \int_{\text{entire ray}} (\varepsilon_\nu/\nu^3)\,\rho_0\,dl
= (1/4\pi)\int (\varepsilon_\nu/\nu^3)_{\text{in region of strongest emission}}
  \rho_0\,dl.
\end{equation}

\begin{equation}
\tag{2.6.38}
\simeq \left(\int \rho_0\,dl\right)\!(1/4\pi)\,
(\varepsilon_\nu/\nu^3)_{\text{in region of strongest emission}}.
\end{equation}

In summary, the radiation from an optically thin source with negligible
scattering $I_\nu \propto \varepsilon_\nu$, and has an intensity proportional
to the amount of matter along the line of sight.  But radiation from an
optically thick source in thermodynamic equilibrium with negligible scattering
has the blackbody form independent of the nature of its emissivity.  Media with
non-negligible scattering will be studied in~\S5.10.

%%% ---------------------------------------------------------------
%%% Radiative transfer in the diffusion approximation
%%% ---------------------------------------------------------------

\bigskip
\noindent
\emph{Radiative transfer in the diffusion approximation}

\medskip
\noindent
Consider the interior of an optically thick medium $\tau_\nu \gg 1$ which is
in local thermodynamic equilibrium.  Suppose that the medium has a temperature
gradient and/or an acceleration, but that the characteristic length scale
%
\begin{equation}
\tag{2.6.39}
l_T \equiv \frac{T}{|\boldsymbol{\nabla}T + \dot{\mathbf{a}}\,T|}
\end{equation}
%
over which the temperature changes are long compared to the mean-free path of a
photon,
%
\begin{equation}
\tag{2.6.40}
l_T \gg 1/\kappa_\nu\,\rho_0 \equiv l_\nu.
\end{equation}
%
Then the radiation distribution will consist of a large, isotropic, blackbody
component, plus a tiny ``correction'' due to the temperature gradient
%
\begin{equation}
\tag{2.6.41}
I_\nu \simeq B_\nu + I_\nu^{(1)};\qquad
I_\nu^{(1)} \ll B_\nu.
\end{equation}

There is no net flux $F$ across any surface associated with the blackbody
component~$B_\nu$; but the ``correction'' term $I_\nu^{(1)}$ will lead to a
flux in the direction of
$\mathbf{h}\cdot(\boldsymbol{\nabla}T+\dot{\mathbf{a}}\,T)$.

(Here $\mathbf{h}$ is the projection operator of~\S2.5.)  One can calculate
the magnitude of that flux by inserting expression~(2.6.41) for~$I_\nu$ into
the equation of transfer~(2.6.22), and by then integrating over
$\cos\theta\,d\Omega\,d\nu$.  (Here $\theta$ is the angle between
$d\Omega$~and the direction of
$\mathbf{h}\cdot(\boldsymbol{\nabla}T+\dot{\mathbf{a}}\,T)$.)
The result is
%
\begin{equation}
\tag{2.6.43}
\mathbf{q} = (1/\kappa_R\rho_0)(\tfrac{1}{3}\partial T^4/\partial \mathbf{x})
  \,\mathbf{h}\cdot(\boldsymbol{\nabla}T + \dot{\mathbf{a}}\,T).
\end{equation}

Here $\mathbf{q}$ is the energy flux vector of~\S2.5, and its magnitude is the
flux of energy in the
$\mathbf{h}\cdot(\boldsymbol{\nabla}T+\dot{\mathbf{a}}\,T)$ direction
%
\begin{equation}
\tag{2.6.44}
|\mathbf{q}| = F.
\end{equation}

Also, in eq.~(2.6.43) $\bar{\kappa}$ is the ``Rosseland mean opacity'',
defined by
%
\begin{equation}
\tag{2.6.45}
\frac{1}{\bar{\kappa}}
\equiv \frac{\displaystyle\int_0^\infty
  (\kappa_\nu + \kappa_s)^{-1}\,
  (dB_\nu/dT)\,d\nu}
  {\displaystyle\int_0^\infty (dB_\nu/dT)\,d\nu}.
\end{equation}

For free-free transitions by themselves (eq.~2.6.25) the Rosseland mean
opacity is
%
\begin{equation}
\tag{2.6.46}
\kappa_{ff} = (0.645\times 10^{23}\;\text{cm}^2/\text{g})
  \left(\frac{C_f\,\bar{f}\,\bar{g}^2}{\mathcal{A}}\right)
  \left(\frac{\rho_0}{\text{g\;cm}^{-3}}\right) T^{-7/2}.
\end{equation}

See~\S2.1 for notation.  For electron scattering by itself in the nonrelativistic
case, the Rosseland mean opacity is
%
\begin{equation}
\tag{2.6.47}
\bar{\kappa}_{es} = \kappa_{es} = 0.40\;\text{cm}^2/\text{g}.
\end{equation}


%%% ===============================================================
%%% SECTION 2.7: Shock Waves
%%% ===============================================================

\subsection{Shock Waves}
\label{sec:2.7}

\emph{Basic references}\quad
Zel'dovich and Raizer~(1966)~\cite{ZeldovichRaizer1966};
Landau and Lifshitz~(1959)~\cite{LandauLifshitz1959};
Taub~(1948)~\cite{Taub1948};
Lichnerowicz~(1967, 1970, 1971)~\cite{Lichnerowicz1967};
Thorne~(1973a)~\cite{Thorne1973a}.

\bigskip
\noindent
\emph{Types of shock waves}

\medskip
\noindent
When gas in supersonic flow encounters an obstacle (e.g., the surface of a star,
or the geometric structure of the ergosphere of a black hole), a shock front
develops.  In the shock front the flow decelerates sharply from supersonic to
subsonic, and some of the kinetic energy of the flow gets converted into heat
(increase in entropy).

The structure of the shock depends on the nature of the forces which decelerate
the gas particles (atoms, ions, electrons).  If those forces are collisions between
the particles themselves (the usual case), one has an ordinary shock.  But if the
particles are decelerated without colliding---e.g., by impact onto the dipole
magnetic field of a neutron star, which swings the particles into Larmour
orbits---one has a \emph{collisionless} shock.  These notes will be confined to ordinary
shocks.  For the theory of collisionless shocks see, e.g., the end of~\S12 of Kaplan
(1966).

Shock waves in a partially ionized gas (e.g., the interstellar medium) can have
a somewhat different form than ordinary shocks.  As gas passes through the shock
front, much of its kinetic energy of supersonic flow can be converted into
excitation energy of atoms and ions, and can be quickly radiated away as ``line''
radiation.  The result is a bright glow from the shock front itself.  See, e.g.,
Pikel'ner~(1961) and~\S10 of Kaplan~(1966).  In this section we shall ignore such
energy losses to radiation---i.e., we shall demand that the gas behind the shock
have the same total energy per unit rest mass as the gas in front of the shock.

%%% ---------------------------------------------------------------
%%% The relativistic Rankine--Hugoniot equations
%%% ---------------------------------------------------------------

\bigskip
\noindent
\emph{The relativistic Rankine--Hugoniot equations}

\medskip
\noindent
Pick a particular event~$\mathscr{P}$ on a shock front.  In the neighborhood
of~$\mathscr{P}$ introduce a local Lorentz frame (``rest frame of the shock'')
in which (i) the shock is the surface $y = z = 0$; and (iii) on both sides of
the shock the fluid is moving in the $x$~direction, i.e., perpendicular to the
shock front (``normal shock''; see Figure~\ref{fig:2.7.1}).  That such a local
Lorentz frame exists in general one can prove quite easily [see, e.g.,
Taub~(1948)].

Denote the ``front'' side of the shock (side \emph{from} which the fluid moves)
by a ``1'', and denote the ``back'' side (side \emph{toward} which the fluid
moves) by a ``2''; see Figure~\ref{fig:2.7.1}.  Denote the velocity of the
fluid, as measured in the rest frame of the shock, by
%
\begin{subequations}
\begin{align}
\tag{2.7.1a}
v_1 &= (dx/dt)_1 = \text{ordinary velocity on front side},\\[4pt]
\tag{2.7.1b}
v_2 &= (dx/dt)_2 = \text{ordinary velocity on back side}.
\end{align}
\end{subequations}
%
(Note that these ``4-velocities'' are scalars, not vectors.)

In the rest frame of the shock the law of baryon conservation is equivalent
to continuity of the \emph{baryon flux}:
%
\begin{equation}
\tag{2.7.2a}
j \equiv n_1\,u_1 = n_2\,u_2,
\end{equation}
%
and the relations $n = 1/V_1$, $\mu_1 = p_1$,
%
\begin{subequations}
\begin{align}
\tag{2.7.1c}
u_1 &= v_1\,\gamma_1 = v_1/(1 - v_1^2)^{1/2}
  = \text{``4-velocity'' on front side},\\[4pt]
\tag{2.7.1d}
u_2 &= v_2\,\gamma_2 = v_2/(1 - v_2^2)^{1/2}
  = \text{``4-velocity'' on back side}.
\end{align}
\end{subequations}

\noindent
Similarly, energy and momentum conservation are equivalent to continuity of
energy flux and momentum flux:
%
\begin{subequations}
\begin{align}
\tag{2.7.2b}
\bigl[(\rho + p)\,v\,u\bigr]_1
  &= \bigl[(\rho + p)\,v\,u\bigr]_2,\\[4pt]
\tag{2.7.2c}
\bigl[(\rho + p)\,u^2 + p\bigr]_1
  &= \bigl[(\rho + p)\,u^2 + p\bigr]_2.
\end{align}
\end{subequations}

%%% ---------------------------------------------------------------
%%% Figure 2.7.1
%%% ---------------------------------------------------------------

\begin{figure}[htbp]
\centering
% \includegraphics[width=\columnwidth]{fig_2_7_1}
\fbox{\parbox{0.9\columnwidth}{\centering [Figure 2.7.1 placeholder]\\[6pt]
Diagram of a shock wave viewed in the local-Lorentz rest frame of the
shock front at an event~$\mathscr{P}$.  The front (``1'' side) is on the left,
and the back (``2'' side) is on the right.  Arrows indicate gas flow with
velocity~$v_1$ entering from the left and velocity~$v_2$ exiting to the right.
The $x$-axis is perpendicular to the shock surface, which lies in the
$y$-$z$ plane.  Hatching indicates the compressed gas behind the shock.}}
\caption{\textit{Figure~2.7.1.}  A shock wave viewed in the local-Lorentz
rest frame of the shock front at an event~$\mathscr{P}$.}
\label{fig:2.7.1}
\end{figure}

These junction conditions are more easily understood by writing them in a
form analogous to the Rankine--Hugoniot equations of Newtonian theory:
First take the law of baryon conservation~(2.7.2a) and turn it into equations
for the fluid 4-velocities
%
\begin{equation}
\tag{2.7.3a}
u_1 = jV_1, \qquad u_2 = jV_2.
\end{equation}
%
(See~\S2.5 for notation.)  Then take the law of momentum conservation~(2.7.2a)
and rewrite it in terms of $\mu$, $V_1$, and $p$ using the law of baryon
conservation~(2.7.2a) to obtain
%
\begin{subequations}
\begin{align}
\tag{2.7.3b}
j^2 &= -\frac{p_2 - p_1}{\mu_2\,V_2 - \mu_1\,V_1},\\[6pt]
\end{align}
\end{subequations}
%
and the relations
$j^2 = (p_2-p_1)/(\mu_2 V_2 - \mu_1 V_1)$.

Finally, use the law of baryon conservation~(2.7.2b) to rewrite the energy
equation~(2.7.3b) by $(\mu_1 V_1 + \mu_2 V_2)$, and combine with
$j = u_1/V_1 = u_2/V_2$; divide equation~(2.7.3b) by the square of the energy
equation $(\mu_2\gamma_2)^2 - (\mu_1\gamma_1)^2 = (p_1 - p_2)(\mu_1 V_1 + \mu_2 V_2)$;
then subtract this from [eq.~(2.7.2a)] to obtain
%
\begin{equation}
\tag{2.7.3c}
\mu_2^2 - \mu_1^2 = (p_2 - p_1)(\mu_1\,V_1 + \mu_2\,V_2).
\end{equation}

\emph{Equations~(2.7.3) are Taub's~(1948) junction conditions for shock waves.}
Their Newtonian limits are the standard Rankine--Hugoniot equations:
%
\begin{subequations}
\begin{align}
\tag{2.7.4a}
v_1 &= j_0\,V_{01},\qquad v_2 = j_0\,V_{02}\\[4pt]
\tag{2.7.4b}
j_0 &\equiv \frac{p_2 - p_1}{V_{01} - V_{02}}\quad
\text{(``mass flux'' of Newtonian theory)};\\[4pt]
\tag{2.7.4c}
w_2 - w_1 &= \tfrac{1}{2}(p_2 - p_1)(V_{01} + V_{02}).
\end{align}
\end{subequations}
%
The form of the Newtonian junction conditions~(2.7.4) motivates one to use $p$
and $V_0$ as one's independent thermodynamic variables when analyzing Newtonian
shocks; similarly, the form of the relativistic junction conditions~(2.7.3)
motivates one to use $p$ and~$\mu V$ as one's independent variables for
relativistic shocks.
